\section{Introduction}\label{sec:intro}

The canonical bundle formula separates the canonical geometry of a
fibration into a discriminant, which measures singularities over divisors,
and a moduli part, which measures variation.
The b-semiampleness conjecture predicts that a suitable multiple of the
moduli part defines a morphism on a birational model of the base.
Hodge-theoretic semipositivity underlies this formula
\cite{Kawamata81}. Ambro established birational stabilization and
positivity of the moduli part in the generic-klt setting
\cite{Ambro04,Ambro05}; Fujino--Gongyo proved b-nefness and abundance
for projective lc-trivial fibrations \cite{FujinoGongyo14}.
In the algebraic category, the conjecture is proved by
Bakker--Filipazzi--Mauri--Tsimerman \cite[Theorem~1.5]{BFMT25}.
Their discussion of the analytic category retains projectivity of the
morphism; the extension to K\"ahler morphisms is explicitly distinguished
in \cite[Remarks~6.31 and~7.5]{BFMT25}.

We resolve the compact log-smooth case of this K\"ahler b-semiampleness
question. Both the source and the base may be nonprojective.
The boundary may have coefficient one, so the relevant Hodge structure
can be mixed. The assertion concerns actual holomorphic line bundles:
numerical semipositivity alone would not suffice.

\subsection{Statement}

All spaces and maps are complex analytic unless stated otherwise.
A \emph{modification} is a proper bimeromorphic holomorphic map.
Write \(\Pic(X)_{\Q}=\Pic(X)\otimes_{\Z}\Q\), using additive notation.
Equality in this group means an isomorphism of holomorphic line bundles
after a common positive integral multiple.

Let \(f:Y\to X\) be a surjective holomorphic map with connected fibers,
where \(Y\) and \(X\) are smooth compact connected K\"ahler manifolds.
Let \(\Delta\) be an effective rational divisor with simple normal
crossing support and coefficients in \([0,1]\), and assume that
\begin{equation}\label{eq:adjoint}
 K_Y+\Delta\sim_{\Q}f^*L,\qquad L\in\Pic(X)_{\Q}.
\end{equation}
These assumptions give an lc-trivial fibration in the usual rank-one
sense; the rank condition is recalled in
Section~\ref{sec:prelim}.

Here is the moduli object in the statement.
For a smooth compact K\"ahler modification \(\mu:X'\to X\), resolve
the main component of \(Y\times_X X'\):
\[
\begin{tikzcd}[column sep=large,row sep=large]
 Y' \arrow[r,"h"] \arrow[d,"f'"'] & Y\arrow[d,"f"]\\
 X' \arrow[r,"\mu"'] & X .
\end{tikzcd}
\]
Take \(Y'\) smooth compact K\"ahler and set
\[
 \Delta'=h^*\Delta-K_{Y'/Y}.
\]
Thus \(K_{Y'}+\Delta'=h^*(K_Y+\Delta)\) under the natural
canonical-bundle identification. Exceptional coefficients of
\(\Delta'\) can be negative.

For a prime divisor \(P\subset X'\), put
\begin{equation}\label{eq:threshold}
 t_P=\inf_{E\mapsto P}
 \frac{a(E;Y',\Delta')}{\ord_E((f')^*P)}.
\end{equation}
Here the infimum ranges over all prime divisors on smooth proper
bimeromorphic models of \(Y'\) that dominate \(P\), with positive
denominator; \(a(E;Y',\Delta')\) denotes log discrepancy.
Equivalently, \(t_P\) is the largest \(t\) for which
\((Y',\Delta'+t(f')^*P)\) is sub-log-canonical over a general point of
\(P\). Define
\begin{equation}\label{eq:moduli}
 B_{X'}=\sum_P(1-t_P)P,\qquad
 M_{X'}=\mu^*L-K_{X'}-B_{X'}\ \in\Pic(X')_{\Q}.
\end{equation}
These definitions are independent of the chosen resolution of the
source. The coefficients are rational and only finitely many are
nonzero; the elementary local calculation is recalled below.
The family \(\mathbf M=(M_{X'})\) is the \emph{moduli rational
b-line bundle}.

\Needspace{8\baselineskip}
\begin{theorem}\label{thm:main}
Under the assumptions above, \(\mathbf M\) is b-semiample.
More explicitly, there is a smooth compact K\"ahler modification
\(\mu_0:S\to X\) such that:
\begin{enumerate}[label=\textup{(\alph*)}]
\item for every smooth compact K\"ahler modification \(\nu:S_1\to S\),
\[
 M_{S_1}=\nu^*M_S\quad\text{in }\Pic(S_1)_{\Q};
\]
\item for some positive integer \(m\), the actual line bundle
representing \(mM_S\) is generated by its global holomorphic sections.
\end{enumerate}
The theorem includes a point base and relative dimension zero.
No projectivity of \(f\), \(X\), or \(Y\) is assumed.
\end{theorem}

The discriminant in \eqref{eq:moduli} is the log-canonical-threshold
discriminant. It is not the weighted inf-multiplicity divisor of an
orbifold base. In particular, a Hodge line produced by a decomposition
of pluriforms must still be identified with \eqref{eq:moduli}.
That identification is one of the principal steps of the proof.

\subsection{Proof and technical contributions}

Choose a compact K\"ahler modification \(S\) whose boundary is simple
normal crossing. Locally on \(S\), take a cyclic root of a generator of
the relative log-plurivolume line. Its highest Hodge eigenspace is
one-dimensional. If horizontal coefficient-one components are present,
an iterated residue identifies that line, through a single pure weight
grade, with the volume line of a deepest boundary stratum.
The mixed Hodge extension theorem of Fujino--Fujisawa
\cite[Theorem~1.1]{FF25} makes this reduction compatible with extension
over a disk.

The local extension calculation keeps track of all divisorial
valuations and of the base Jacobian. For a local root volume \(\tau\),
it proves the precise normalized order
\[
 \ord_P^{\mathrm{Hodge}}(\tau)=t_P-1.
\]
It follows that the canonical Hodge extension is the actual moduli
line and that it pulls back on every further smooth model. The same
calculation works with arbitrary crepant vertical coefficients.
It is therefore also available later on a projective comparison family.

Adjunction to a deepest stratum produces smooth compact K\"ahler klt
fibers with trivial log-pluricanonical bundle.
A theorem of Matsumura--Wang--Wu--Zhang gives finite product covers of
these fibers \cite[version~1, Corollary~1.3]{MWWWZ25}.
A pointwise product decomposition is insufficient for semiampleness.
We construct a marked family with a section, realize its finite fiber
covers after finite base changes, and parameterize their product graphs
inside fixed compact K\"ahler ambients.
Compact Douady components and a proper-constructibility argument give
a compact proper surjection \(P\to S\) carrying the required product
family on a dense open. The auxiliary base need not be projective or
generically finite over \(S\).

The product has a projective log Calabi--Yau factor and factors whose
highest forms are controlled by weight-one or weight-two periods.
The first is handled by the algebraic b-semiampleness theorem on a
projective comparison base. For the others, a flat change of the Hodge
filtration turns a real polarization into a rational one without
changing the underlying integral monodromy. Arithmetic period
compactifications then supply semiample extensions. This use of
period compactification is analytic on the auxiliary base
\cite[Theorems~5.2 and~5.5]{BFMT25}.

Finally, a product identity of volume norms and two-sided logarithmic
bounds extend a specified open isomorphism of line bundles across the
boundary. This step excludes an undetected flat twist.
Semiampleness descends from \(P\) through Stein factorization and finite
analytic norms. These arguments yield global generation at every point
of \(S\), with one exponent.
The proof does not use the Campana orbifold Iitaka theorem.
